\documentclass[../main.tex]{subfiles}
\begin{document}
This appendix provides detailed use case specifications for the core functions of the Smart Medical QA \& Consultation System.

\section{Use Case UC-01: Consult Medical Agent}
\label{section:uc-01}
Use Case UC‑01 – Consult Medical Agent: Actors include clinicians, patients, or administrative staff. The user queries the intelligent agent in natural language, triggering a localized ReAct loop that searches the custom Neo4j Knowledge Graph, retrieves clinical evidence chunks, inspects pill details, and constructs a structured Vietnamese response with Vis.js network citations. Preconditions are that the local Neo4j service is running with a KG of 12,942 entities and the local backend hosts the 3-billion parameter model. In the normal flow, the user enters a clinical question (e.g., "Diabetes diet suggestions"), the system detects intent and invokes the ReAct pipeline, the agent generates a thought and runs the \texttt{graphrag\_query} tool, Cypher queries retrieve matching entities and context paragraphs, the agent synthesizes evidence in a second iteration and outputs the \texttt{Final Answer}, and the Web UI streams markdown, renders a Vis.js subgraph, and offers PDF export. Alternative flows include format recovery (repairing missing ReAct tags), loop‑guard trigger (terminating duplicate tool calls), and intent routing fallback (direct quantitative lab checks).

\section{Use Case UC-02: Import and Analyze Multi-format Lab Report}
\label{section:uc-02}
Use Case UC‑02 – Import and Analyze Multi‑format Lab Report: Actors are clinicians or patients. The user uploads a digital lab report in PDF or Excel format; the system extracts textual metrics, automatically compares them against on‑form reference ranges, highlights abnormal parameters, and generates a personalized clinical consultation using graph‑derived recommendations. The precondition is a valid local file containing standard laboratory grids. In the normal flow, the user drops the file onto the Web UI, the FastAPI server streams it to the PyMuPDF or openpyxl extractor, the extracted grids are processed by the \texttt{lab\_compare\_on\_form} module, the system flags indicators as normal, high, or low, extracts abnormal entities, queries Neo4j for recommendations, prompts the LLM for an advice plan, and the UI displays side‑by‑side color‑coded alerts and advisory text.

\section{Use Case UC-03: Manage Medication Schedule and Reminder}
\label{section:uc-03}
Use Case UC‑03 – Manage Medication Schedule and Reminder: Actors are patients or caregivers. The user sets up daily intake schedules for suggested medications, visualizes schedule blocks on a weekly calendar, and receives automated audio‑visual reminders. The precondition is that the agent has generated a prescription plan or the user manually inputs drug names. In the normal flow, the user opens the Reminder Calendar, selects a medication card, defines frequency, target times, and intake markers, the system stores the configuration in browser local storage, the Web UI updates the weekly schedule with colored blocks, and the browser triggers an alert badge and audio tone when a time slot is reached.

\section{Use Case UC-04: Clinical Calculators and Interaction Safety Check}
\label{section:uc-04}
Use Case UC‑04 – Clinical Calculators and Interaction Safety Check: Actors are clinicians or patients. The user employs deterministic tools to calculate standard metrics (BMI, eGFR) or queries the Neo4j Knowledge Graph to check drug‑drug interaction safety. Preconditions are that the user has accessed the "Công cụ" page and the backend API is connected to Neo4j. In the normal flow, the user navigates to the clinical tools interface, inputs weight and height to calculate BMI (instant result with classification), inputs serum creatinine, age, and sex to compute eGFR using the CKD‑EPI (2021) equation, or selects two drug names to check interaction safety, where the backend queries Neo4j and returns clinical warning text and source evidence.

\end{document}